\begin{abstract}
O monitoramento da qualidade do ar constitui um componente central da gestão ambiental urbana contemporânea, em razão da relação direta entre a poluição atmosférica e os impactos adversos à saúde pública. Neste trabalho, propõe-se uma abordagem multivariada para a classificação da qualidade do ar, baseada na construção de um índice de poluição composto e em sua posterior discretização em classes ordinais. O estudo utiliza o conjunto de dados Air Quality do repositório UCI Machine Learning Repository, que reúne medições horárias de poluentes de referência, respostas de sensores químicos de óxido metálico (MOx) e variáveis meteorológicas coletadas em um ambiente urbano. Um índice de poluição robusto é construído a partir das concentrações dos poluentes de referência por meio de normalização baseada na mediana e no intervalo interquartil, sendo posteriormente discretizado em cinco níveis de poluição por quantis empíricos. A tarefa de classificação é abordada por modelos supervisionados lineares e não lineares, incluindo Regressão Logística Multinomial, Análise Discriminante Linear e Quadrática, k-Nearest Neighbors e Multi-Layer Perceptron. O desempenho dos modelos é avaliado em um conjunto de teste independente utilizando acurácia, F1-score por classe e F1-macro. Os resultados indicam desempenho consistente entre os modelos, especialmente para os níveis intermediários de poluição, com o Multi-Layer Perceptron apresentando a melhor acurácia global e F1-macro. A análise das matrizes de confusão revela que os erros de classificação ocorrem predominantemente entre classes adjacentes, preservando a estrutura ordinal dos níveis de poluição. De forma geral, os resultados demonstram que o índice composto proposto constitui uma representação eficaz da poluição atmosférica global e uma base confiável para sistemas de classificação multiclasse da qualidade do ar em ambientes urbanos.
\end{abstract}


\begin{IEEEkeywords}
qualidade do ar; poluição atmosférica; sensores ambientais; índice de poluição composto; classificação multiclasse; aprendizado de máquina; regressão logística; análise discriminante linear; k-nearest neighbors; multilayer perceptron; análise discriminante quadrática.
\end{IEEEkeywords}
